% !TeX root = ../informe.tex
% ===========================================================================
%  Anexos
% ===========================================================================

\section{Scripts de base de datos}
\label{anexo:sql}

Los scripts completos se mantienen versionados junto con la implementación.
Las siguientes rutas permiten consultar la versión exacta aplicada por cada
microservicio, incluidos los datos semilla:

\begin{itemize}
  \item Inicialización de PostgreSQL:
  \url{https://github.com/miguesj-hub/reserva-canchas/tree/main/infra/postgres/init}
  \item Migraciones de usuarios:
  \url{https://github.com/miguesj-hub/reserva-canchas/tree/main/backend/ms-usuarios/src/main/resources/db/migration}
  \item Migraciones de canchas:
  \url{https://github.com/miguesj-hub/reserva-canchas/tree/main/backend/ms-canchas/src/main/resources/db/migration}
  \item Migraciones de reservas:
  \url{https://github.com/miguesj-hub/reserva-canchas/tree/main/backend/ms-reservas/src/main/resources/db/migration}
\end{itemize}

\section{Colección de pruebas de la API}
\label{anexo:api}

La colección Postman contiene las peticiones de prueba que atraviesan el API
Gateway, incluidas las variables de conexión, autenticación y los endpoints
de usuarios, canchas, reservas y reportes:

\url{https://github.com/miguesj-hub/reserva-canchas/blob/main/Postman-collection/reserva-canchas.postman_collection.json}

\section{Manual de despliegue}
\label{anexo:despliegue}

Este anexo resume el procedimiento de puesta en marcha del sistema, de modo que
el informe sea autosuficiente para reproducir el entorno descrito en los
capítulos anteriores. El manual completo, con la guía de operación y el
diagnóstico de fallos, constituye el entregable \ent{5} y se encuentra en
\id{despliegue/despliegue.pdf} dentro del repositorio.

\subsection{Requisito previo}
\label{anexo:despliegue:requisito}

El único requisito es \textbf{Docker 24 o superior con Compose v2}. No es
necesario instalar Java, Node, Maven ni PostgreSQL en la máquina anfitriona:
la totalidad del sistema se compila y se ejecuta dentro de contenedores. Esta
decisión responde al principio de arranque con un solo comando adoptado por el
proyecto, y evita que la reproducción del entorno dependa de las herramientas
instaladas en cada equipo.

\subsection{Puesta en marcha}
\label{anexo:despliegue:arranque}

\begin{lstlisting}[
  numbers=none,
  caption={Clonado y arranque del sistema completo.},
  label={lst:despliegue-arranque}
]
git clone https://github.com/miguesj-hub/reserva-canchas
cd reserva-canchas
docker compose up -d --wait
\end{lstlisting}

La aplicación queda disponible en \url{http://localhost}. El primer arranque
compila las imágenes y requiere varios minutos; los posteriores reutilizan las
capas ya construidas. Al finalizar, \id{docker compose ps} debe mostrar los
once servicios en estado \id{healthy}.

Todos los puertos publicados se pueden reasignar mediante variables de entorno,
sin modificar el archivo de Compose. Por ejemplo, si el puerto 8080 ya está
ocupado en la máquina anfitriona:

\begin{lstlisting}[
  numbers=none,
  caption={Reasignación del puerto del \emph{gateway}.},
  label={lst:despliegue-puerto}
]
GATEWAY_PORT=18080 docker compose up -d --wait
\end{lstlisting}

Esta reasignación no afecta al acceso a la aplicación: el navegador utiliza
únicamente \url{http://localhost}, que corresponde al puerto 80 del contenedor
\id{edge}. Los puertos de los microservicios se publican solamente para poder
inspeccionar su documentación Swagger.

\subsection{Servicios desplegados}
\label{anexo:despliegue:servicios}

La composición levanta once contenedores. Nueve se construyen desde el propio
repositorio y dos se descargan ya construidos, según se detalla en la
\cref{tab:anexo-contenedores}.

\begin{table}[H]
\centering
\caption{Contenedores que componen el despliegue.}
\label{tab:anexo-contenedores}
\small
\begin{tabularx}{\textwidth}{@{}l L{2.1cm} c Y@{}}
\toprule
\textbf{Servicio} & \textbf{Red} & \textbf{Puerto} & \textbf{Responsabilidad} \\
\midrule
\id{edge}              & frontend y backend & 80   & Origen único: sirve el \emph{shell}, los \emph{remotes} y \id{/api} \\
\id{shell}             & frontend & 3000 & Anfitrión de los microfrontends, sesión y enrutado \\
\id{mf-reservas}       & frontend & 3001 & Disponibilidad, creación y consulta de reservas \\
\id{mf-administracion} & frontend & 3002 & Canchas, reservas y usuarios \\
\id{mf-reportes}       & frontend & 3003 & Indicadores de ocupación y uso \\
\id{gateway}           & backend  & 8080 & Verificación de credenciales y enrutado \\
\id{ms-usuarios}       & backend  & 8081 & Registro, autenticación y gestión de usuarios \\
\id{ms-canchas}        & backend  & 8082 & Catálogo, horarios y bloqueos \\
\id{ms-reservas}       & backend  & 8083 & Reservas y disponibilidad \\
\id{ms-reportes}       & backend  & 8084 & Indicadores; sin base de datos propia \\
\id{postgres}          & backend  & 5432 & Las tres bases, aisladas por credenciales \\
\bottomrule
\end{tabularx}
\end{table}

El contenedor \id{edge} es el único conectado a ambas redes. Esa disposición
impide que el navegador alcance un microservicio sin atravesar el
\emph{gateway}, tal como se describe en el capítulo de arquitectura.

\subsection{Usuarios de prueba}
\label{anexo:despliegue:usuarios}

La migración \id{V2__seed_usuarios.sql} siembra las cuentas de la
\cref{tab:anexo-usuarios} durante el primer arranque. Las contraseñas se
almacenan cifradas con BCrypt; en la base de datos no se conserva ninguna en
claro. Se trata de credenciales de demostración, destinadas exclusivamente a la
evaluación del proyecto.

\begin{table}[H]
\centering
\caption{Cuentas sembradas para la evaluación.}
\label{tab:anexo-usuarios}
\small
\begin{tabularx}{\textwidth}{@{}l l L{4.0cm} Y@{}}
\toprule
\textbf{Usuario} & \textbf{Contraseña} & \textbf{Rol} & \textbf{Propósito} \\
\midrule
\id{admin}    & \id{admin}    & Administrador & Catálogo, reservas del club, usuarios y reportes \\
\id{cliente1} & \id{cliente1} & Usuario final & Recorrido de reserva y cancelación \\
\id{cliente2} & \id{cliente2} & Usuario final & Segundo actor del escenario de concurrencia \\
\id{inactivo} & \id{inactivo} & Usuario final, inactivo & Rechazo de acceso de una cuenta inactiva \\
\bottomrule
\end{tabularx}
\end{table}

El catálogo de canchas y un conjunto de reservas de ejemplo también se siembran.
Las reservas emplean fechas relativas al momento del arranque, de manera que las
confirmadas permanezcan en el futuro y puedan cancelarse durante la
demostración.

\subsection{Verificación del despliegue}
\label{anexo:despliegue:verificacion}

\begin{lstlisting}[
  numbers=none,
  caption={Comprobaciones posteriores al arranque.},
  label={lst:despliegue-verificacion}
]
docker compose ps
curl -i http://localhost/api/canchas
curl -u cliente1:cliente1 http://localhost/api/canchas
\end{lstlisting}

La primera orden debe mostrar los once servicios sanos. La segunda debe
responder \id{401}, lo que confirma que el \emph{gateway} rechaza las peticiones
sin credenciales. La tercera devuelve el catálogo sembrado y constituye la
comprobación más informativa del conjunto: si responde correctamente, entonces
el \id{edge}, el \emph{gateway}, \id{ms-usuarios} —que validó la credencial—,
\id{ms-canchas} y PostgreSQL se encuentran operativos y comunicados entre sí.

\subsection{Operación habitual}
\label{anexo:despliegue:operacion}

\begin{lstlisting}[
  numbers=none,
  caption={Órdenes de uso frecuente durante la operación.},
  label={lst:despliegue-operacion}
]
docker compose down -v && docker compose up -d --wait
docker compose down
docker compose logs -f ms-reservas
docker compose up -d --wait --build ms-reservas
\end{lstlisting}

La primera orden reinicia el sistema desde cero: el modificador \id{-v} elimina
el volumen de datos, con lo cual vuelven a ejecutarse los guiones de
inicialización y las migraciones con su carga semilla. Es el procedimiento
recomendado antes de una demostración, puesto que las reservas de ejemplo se
generan con fechas relativas al arranque. La segunda detiene el sistema
conservando los datos, la tercera consulta el registro de un servicio y la
cuarta reconstruye la imagen de un servicio tras modificar su código; omitir
\id{--build} en ese último caso es la causa más frecuente de que un cambio no
se refleje en la ejecución.

\subsection{Diagnóstico de fallos frecuentes}
\label{anexo:despliegue:problemas}

\begin{table}[H]
\centering
\caption{Fallos frecuentes durante el despliegue y su tratamiento.}
\label{tab:anexo-problemas}
\small
\begin{tabularx}{\textwidth}{@{}L{4.1cm} L{4.1cm} Y@{}}
\toprule
\textbf{Síntoma} & \textbf{Causa} & \textbf{Tratamiento} \\
\midrule
\id{port is already allocated} & Otro proceso ocupa el puerto & Reasignarlo mediante la variable de entorno correspondiente \\
Respuesta \id{502} en \id{/api} & El \emph{gateway} aún no está sano & Esperar; el \id{edge} resuelve el destino en cada petición y se recupera solo \\
Respuesta \id{503} en \id{/api} & El microservicio destino no responde & Consultar el registro del servicio \\
Un microservicio reinicia de forma reiterada & Falló una migración o la base rechazó la conexión & Consultar su registro; si se modificaron contraseñas, reiniciar con \id{down -v} \\
Un cambio de código no se refleja & Se levantó el servicio sin \id{--build} & Reconstruir la imagen del servicio \\
La sesión se cierra de forma inesperada & La cuenta fue desactivada & El \emph{gateway} valida la credencial en cada petición; volver a iniciar sesión \\
\bottomrule
\end{tabularx}
\end{table}

\section{Capturas de la aplicación}
\label{anexo:capturas}

Las siguientes capturas evidencian, sobre el sistema desplegado con
\texttt{docker compose up}, los flujos principales descritos en el
\emph{quickstart} del proyecto: el escenario V1 (reservar y cancelar) desde
la sesión de un usuario final, y los escenarios V3, V4 y V5 (catálogo,
reportes y roles) desde la sesión del administrador.

\subsection{Usuario final}
\label{anexo:capturas:usuario}

\begin{figure}[H]
  \centering
  \includegraphics[width=0.7\linewidth]{capturas/usuario-login}
  \caption{Inicio de sesión como usuario final (\texttt{cliente1}).}
  \label{fig:cap-usuario-login}
\end{figure}

\begin{figure}[H]
  \centering
  \includegraphics[width=\linewidth]{capturas/usuario-disponibilidad}
  \caption{Disponibilidad de una cancha: bloques libres, ocupados y en
           mantenimiento sobre el horario de atención.}
  \label{fig:cap-usuario-disponibilidad}
\end{figure}

\begin{figure}[H]
  \centering
  \includegraphics[width=\linewidth]{capturas/usuario-seleccion-bloque}
  \caption{Selección de un bloque libre para reservar.}
  \label{fig:cap-usuario-seleccion}
\end{figure}

\begin{figure}[H]
  \centering
  \includegraphics[width=0.75\linewidth]{capturas/usuario-confirmar-reserva}
  \caption{Confirmación de los datos de la reserva antes de crearla.}
  \label{fig:cap-usuario-confirmar}
\end{figure}

\begin{figure}[H]
  \centering
  \includegraphics[width=\linewidth]{capturas/usuario-reserva-confirmada}
  \caption{Reserva creada: aparece como \textsc{Confirmada} en Mis Reservas,
           con la acción de cancelar disponible.}
  \label{fig:cap-usuario-confirmada}
\end{figure}

\begin{figure}[H]
  \centering
  \includegraphics[width=\linewidth]{capturas/usuario-reserva-cancelada}
  \caption{Reserva cancelada por el propio usuario: pasa a estado
           \textsc{Cancelada}.}
  \label{fig:cap-usuario-cancelada}
\end{figure}

\subsection{Administrador}
\label{anexo:capturas:admin}

\begin{figure}[H]
  \centering
  \includegraphics[width=0.7\linewidth]{capturas/admin-login}
  \caption{Inicio de sesión como administrador.}
  \label{fig:cap-admin-login}
\end{figure}

\begin{figure}[H]
  \centering
  \includegraphics[width=\linewidth]{capturas/admin-dashboard}
  \caption{Panel principal de administración con los indicadores del día.}
  \label{fig:cap-admin-dashboard}
\end{figure}

\begin{figure}[H]
  \centering
  \includegraphics[width=\linewidth]{capturas/admin-canchas-listado}
  \caption{Gestión de canchas: catálogo con su horario y estado.}
  \label{fig:cap-admin-canchas}
\end{figure}

\begin{figure}[H]
  \centering
  \includegraphics[width=0.6\linewidth]{capturas/admin-nueva-cancha}
  \caption{Alta de una nueva cancha, con deporte y horario de atención.}
  \label{fig:cap-admin-nueva-cancha}
\end{figure}

\begin{figure}[H]
  \centering
  \includegraphics[width=\linewidth]{capturas/admin-cancha-creada}
  \caption{La cancha creada aparece de inmediato en el catálogo.}
  \label{fig:cap-admin-cancha-creada}
\end{figure}

\begin{figure}[H]
  \centering
  \includegraphics[width=\linewidth]{capturas/admin-reservas-listado}
  \caption{Gestión de reservas: el administrador puede cancelar cualquier
           reserva del club.}
  \label{fig:cap-admin-reservas}
\end{figure}

\begin{figure}[H]
  \centering
  \includegraphics[width=0.65\linewidth]{capturas/admin-confirmar-cancelacion}
  \caption{Confirmación antes de cancelar una reserva de un usuario final.}
  \label{fig:cap-admin-confirmar-cancelacion}
\end{figure}

\begin{figure}[H]
  \centering
  \includegraphics[width=\linewidth]{capturas/admin-reserva-cancelada}
  \caption{La reserva queda \textsc{Cancelada} y el bloque vuelve a estar
           libre.}
  \label{fig:cap-admin-reserva-cancelada}
\end{figure}

\begin{figure}[H]
  \centering
  \includegraphics[width=\linewidth]{capturas/admin-reportes-indicadores}
  \caption{Reportes: total de reservas, ocupación media, cancelaciones y
           cancha de mayor demanda para el rango seleccionado.}
  \label{fig:cap-admin-reportes-indicadores}
\end{figure}

\begin{figure}[H]
  \centering
  \includegraphics[width=\linewidth]{capturas/admin-reportes-ranking}
  \caption{Ranking de demanda por cancha, con sus extremos de mayor y menor
           ocupación.}
  \label{fig:cap-admin-reportes-ranking}
\end{figure}

\begin{figure}[H]
  \centering
  \includegraphics[width=\linewidth]{capturas/admin-usuarios-listado}
  \caption{Gestión de usuarios: roles y estado de cada cuenta.}
  \label{fig:cap-admin-usuarios}
\end{figure}

\begin{figure}[H]
  \centering
  \includegraphics[width=\linewidth]{capturas/admin-usuario-activado}
  \caption{Activación de la cuenta \texttt{inactivo}, que pasa a poder
           iniciar sesión (FR-005, FR-006).}
  \label{fig:cap-admin-usuario-activado}
\end{figure}

\section{Repositorio del proyecto}
\label{anexo:repositorio}

La implementación, las migraciones de base de datos y la colección de
pruebas Postman se encuentran disponibles en el siguiente repositorio:

\url{https://github.com/miguesj-hub/reserva-canchas}
